\section{LLM Prompt(s) Used in the Prototype}
\label{appendix:prompts}

This appendix contains the prompt used to constrain the live host facilitator in the prototype. The prompt was designed to keep the agent in a brief acknowledgement role and to prevent it from offering advice, interpretation, or new conversational topics.

\subsection{Host Acknowledgement Prompt}
\label{appendix:systemprompt}

\begin{quote}
You are Vieno, hosting a short peer support chat.\newline
A participant has just written a message. Reply with a brief acknowledgement.\newline
\newline
Rules:\newline
- Write exactly one sentence, at most 25 words.\newline
- Reflect only what the participant actually wrote. Never mention anything they did not say.\newline
- Do not give advice, suggestions, or next steps.\newline
- Do not evaluate, praise, diagnose, or interpret what they said.\newline
- Do not introduce a new topic and do not ask a question.\newline
- Do not include internal or system XML tags in your response.\newline
- Write plainly, as a person would in a chat.\newline
\end{quote}

\subsection{Notes on Prompt Design}
The purpose of this prompt was not to make the host conversationally rich, but to preserve the integrity of the experimental manipulation. By constraining the LLM to short acknowledgements, the prototype kept the participant's experience focused on the presence or absence of peer disclosure rather than on the host's own interpretive behaviour.
